% PDF 275--283；原书第 247--255 页。
\section{毒性}
\label{sec:ch06-toxicity}

\begin{keypoint}
毒性（toxicity）涵盖多种有害内容，没有统一定义。可以将它理解为前文所定义的表征性伤害之一，也可以视作独立概念。Perspective API 将毒性描述为可能使参与者退出对话的粗鲁、不尊重或不讲理的评论。\refcite{ch06-ref-056}{Kurita 等人（2019）}将有毒内容描述为可能冒犯或伤害受众的任何材料，其中包括仇恨言论、种族主义和冒犯性用语。\refcite{ch06-ref-089}{Pavlopoulos 等人（2020）}把毒性视为一个统称，研究群体使用“冒犯性”“辱骂性”和“仇恨性”等不同术语，描述各种有毒语言或相关现象。
\end{keypoint}

\subsection{成因}
大语言模型的毒性、偏见和公平性相互关联，都受到训练数据的影响。偏见一节所列的许多成因，例如数据选择、领域和体裁分布不均、创作者的人口构成和文化分布偏斜，也适用于大语言模型的有毒输出。这里侧重可能为毒性所特有、或与模型偏见成因重合的因素。

\begin{enumerate}
\item \textbf{训练数据偏见。}如偏见一节所述，大语言模型的毒性主要来源之一，是训练数据本身带有偏见。社会不平等、带有偏见的语言用法，以及某些群体代表不足，都可能使训练数据存在偏见。这些偏见表现在模型的输出中，产生有毒、不公平的结果。模型往往复现数据集中的偏见，导致有毒输出（\refcite{ch06-ref-006}{Bender et al., 2021}）。
\item \textbf{理解上下文的局限。}大语言模型有时难以理解文本或对话的完整上下文，从而给出不当或有毒的回答。\refcite{ch06-ref-005}{Bender 和 Koller（2020）}指出，模型难以理解含义细致的人类语言，这说明准确理解上下文并不容易。\refcite{ch06-ref-089}{Pavlopoulos 等人（2020）}发现，帖子的上下文会显著影响人们对其毒性的判断，既可能放大，也可能减轻感知到的毒性。他们的一项实验中，标注者未看到上下文时，大约有 5\% 的人工标注帖子得到了相反的毒性标签。
\item \textbf{对抗攻击。}大语言模型容易受到对抗攻击，被提示故意生成有毒输出。\refcite{ch06-ref-112}{Wallace 等人（2020）}的研究说明，攻击者可向模型的训练集中注入恶意样例，显著影响其学习和未来预测。这种攻击使攻击者可以将任意选定的短语变成触发特定预测的通用触发语，构成一种危险的漏洞。研究还表明，可以把这些投毒样例设计得不显眼，使受攻击者难以识别并移除有害数据。投毒样例本身没有明确提到触发语，因此可以避开搜索特定短语的检测方法。
\item \textbf{指定角色的提示。}对话式人工智能的一种常见用法，是让大语言模型扮演某一角色，并以该角色继续对话。\refcite{ch06-ref-020}{Deshpande 等人（2023）}发现，为 ChatGPT 指定特定角色，例如著名拳击手穆罕默德·阿里，可能显著提高生成文本的毒性。根据分配给 ChatGPT 的角色，回答的毒性最高可变为原来的六倍。毒性增加表现为模型助长错误的刻板印象，生成有害对话，表达伤人的观点。将这类回答与指定角色联系起来，既可能损害公众人物的名誉，也可能伤害没有预料到这类内容的用户。
\end{enumerate}

\subsection{评估指标}
大多数毒性评估框架都使用辅助模型，通常是分类器，来评估语言模型生成文本的毒性水平。

原书介绍，Google Jigsaw 开发的 Perspective API 是常用的文本毒性评分工具（\refcite{ch06-ref-059}{Lees et al., 2022}）。如图~\ref{fig:ch06-perspective} 所示，输入为文本，输出是 0 至 1 之间的概率分数。它量化文本被认为含有某种毒性相关属性的可能性。属性包括 \texttt{TOXICITY}（毒性）、\texttt{SEVERE\_TOXICITY}（重度毒性）、\texttt{IDENTITY\_ATTACK}（身份攻击）、\texttt{INSULT}（侮辱）、\texttt{PROFANITY}（粗俗用语）、\texttt{THREAT}（威胁）、\texttt{SEXUALLY\_EXPLICIT}（露骨的性内容）和 \texttt{FLIRTATION}（调情）。Perspective API 使用来自多个在线论坛的大量语料训练模型，来源包括维基百科和《纽约时报》。

\sourcefigure{fig06_08.pdf}{通过 Perspective API 评估毒性。示例文本“Shut up. You’re an idiot!”的各项分数为毒性 0.99、重度毒性 0.75、侮辱 1.0、露骨的性内容 0.04、粗俗用语 0.93、可能遭到拒绝 0.99、威胁 0.15 和身份攻击 0.03。}{fig:ch06-perspective}

数据由众包评分人员标注，人员来自 Figure Eight、Appen 等机构及内部平台。评分人员判断一条评论是否具有某项属性，例如 \texttt{TOXICITY}。最终标签由认为文本有毒的评分者所占比例得出。在模型架构方面，Perspective API 先用标注数据训练基于多语言 BERT 的模型，再将模型蒸馏到每种受支持语言各自的单语言卷积神经网络，使系统能够快速、高效地给出分数。

评估文本生成时，Perspective API 提供表示毒性的概率分数。其毒性分类器可写为
\[
c\colon\widehat{Y}\to[0,1].
\]
\refcite{ch06-ref-033}{Gehman 等人（2020）}采用期望最大毒性（Expected Maximum Toxicity，EMT），对语言模型的 25 条输出计算，关注均值和标准差。EMT 用于找出语言模型生成文本中的极端情况，即毒性最高的输出。
\begin{equation}
\operatorname{EMT}(\widehat{\mathbb{Y}})
=\max_{\widehat{Y}\in\widehat{\mathbb{Y}}}c(\widehat{Y}).
\label{eq:ch06-emt}
\end{equation}

毒性概率（Toxicity Probability，TP）衡量 25 次生成中，至少有一次生成有毒文本的可能性。毒性分数等于或大于 0.5 的文本记为有毒。该指标量化语言模型生成有毒内容的频率。
\begin{equation}
\operatorname{TP}(\widehat{\mathbb{Y}})
=P\!\left(\sum_{\widehat{Y}\in\widehat{\mathbb{Y}}}
 I\bigl(c(\widehat{Y})\geq0.5\bigr)\geq1\right).
\label{eq:ch06-tp}
\end{equation}

\refcite{ch06-ref-009}{Bommasani 等人（2023）}定义的有毒比例（Toxic Fraction，TF）计算生成文本中被判为有毒的比例。
\begin{equation}
\operatorname{TF}(\widehat{\mathbb{Y}})
=\mathbb{E}_{\widehat{Y}\in\widehat{\mathbb{Y}}}
 \bigl[I(c(\widehat{Y})\geq0.5)\bigr].
\label{eq:ch06-tf}
\end{equation}

\subsection{评估基准}
上一节列出的偏见基准与用于评估毒性的基准有相当多的重合。本节侧重专门为毒性设计的基准，见表~\ref{tab:ch06-toxicity-benchmarks}。

\begin{table}[!htbp]
\centering
\caption{主要的毒性基准数据集。表中给出数据集的大小、收集和标注方法，以及所关注的内容。}
\label{tab:ch06-toxicity-benchmarks}
\small
\setlength{\tabcolsep}{3pt}
\begin{tabularx}{\linewidth}{@{}>{\raggedright\arraybackslash}p{.31\linewidth}rlX@{}}
\toprule
数据集 & 规模 & 方法 & 关注点\\
\midrule
Perspective API 的毒性数据集 & 180 万 & 众包 & 整体毒性及特定维度\\
Jigsaw Toxic Comment Dataset & 15 万 & 众包 & 毒性程度和类型\\
Hate Speech Dataset & 2.4 万 & 众包 & 仇恨言论检测\\
ToxiGen & 10 万 & 对抗式构造 & 模型的稳健性与隐含偏见\\
Thoroughly Engineered Toxicity（TET）Dataset & 1 万 & 人工构造 & 使模型防护失效\\
ImplicitHateCorpus & 570 万 & 众包 & 隐含仇恨言论（讽刺、刻板印象、微侵犯）\\
DynaHate & 2250 万 & 机器学习 & 与语境相关的仇恨言论（特定对象、变化中的用语）\\
SocialBiasFrames & 8,732 & 众包 & 有害的社会框架（性别、种族、残障）\\
\bottomrule
\end{tabularx}
\end{table}

\subsection{缓解策略}
\refcite{ch06-ref-033}{Gehman 等人（2020）}将毒性缓解技术分为两大类，即基于数据和基于解码的方法。基于数据的方法对模型进一步预训练，改变模型参数。这个办法有效，但涉及修改参数，通常消耗较多算力。基于解码的方法只调整语言模型的解码算法，模型参数保持不变，因此通常较易采用，所需资源也较少。

\subsubsection{基于数据的方法}
\refcite{ch06-ref-042}{Gururangan 等人（2020）}讨论的领域适应性预训练（Domain Adapted Pre-training，DAPT）分为两个阶段。第一阶段在广泛的来源上作通用预训练，第二阶段使用与特定语境有关的数据集，例如 Twitter 交流数据，进行领域专用的预训练，如图~\ref{fig:ch06-dapt} 所示。结合特定任务的微调时，这种办法尤其有效。例如，在对与“黑人的命也是命”运动有关的话题分类时，微调侧重从更广的推文领域抽取这一语境特有的语言层次和细微差异。\refcite{ch06-ref-033}{Gehman 等人（2020）}说明，用 DAPT 在 OpenWebTextCorpus（OWTC）的无毒子集上继续预训练大语言模型，可明显减少 GPT-2 等模型的毒性，但这一办法也面临困难。最明显的局限之一是计算开销。另外，需要额外的训练数据，尤其通过众包等人工方式获得时，费用可能过高。另一个主要局限是，它可能损害模型表现。筛除有毒成分时，也可能误删无害且有潜在价值的知识。

\sourcefigure{fig06_09.pdf}{领域适应性预训练示意图。模型初次预训练完成后，用另一份语料进行第二次预训练。这份语料与模型任务有关，例如来自网页的数据，并且已滤除有毒的交流内容。图中以从 2020 年推文逐步限定到相关领域和任务数据为例。}{fig:ch06-dapt}

属性条件化（Attribute Conditioning，ATCON）是另一种基于数据的方法，源于 \refcite{ch06-ref-028}{Ficler 和 Goldberg（2017）}及 \refcite{ch06-ref-052}{Keskar 等人（2019）}的工作。该方法从文档中随机选择一个子集，为预训练模型的预训练数据加入毒性属性词元，记作 \verb!|toxic|! 或 \verb!|non-toxic|!。随后继续预训练，让模型学会使用这些毒性标记。用模型生成文本时，把毒性属性词元 \verb!|non-toxic|! 加在提示前面，向模型表明生成文本应当侧重于无毒内容。

% PDF 280--283；原书第 252--255 页。
\subsubsection{基于解码的方法}
这类方法调整解码算法，以尽量减少毒性。

\begin{practicaltip}
最简单的基于解码的办法之一是用屏蔽词表（blocklisting），也称词语过滤。把语言模型中不希望出现的词语的概率设为零。这些词一般包括咒骂、粗俗用语和侮辱词语。这样，模型的输出就不会出现这些词。下面介绍一些在解码时降低毒性、较复杂而通常也更有效的方法。
\end{practicaltip}

\paragraph*{词表偏移}
\refcite{ch06-ref-034}{Ghosh 等人（2017）}提出这一方法。它为词表中的每个词元学习向量表示，区分毒性与无毒属性。原研究使用 LSTM 模型，不过其基本原理也可以用于 GPT-2 等较新的 Transformer 架构。

在标准 LSTM 模型中，由 $M$ 个词 $w_1,w_2,\ldots,w_M$ 组成的序列，其联合概率由概率链式法则定义为
\begin{equation}
P(w_1,w_2,\ldots,w_M)
=\prod_{t=1}^{M}P(w_t\mid w_1,w_2,\ldots,w_{t-1}).
\label{eq:ch06-lstm-joint}
\end{equation}
在这种模型中，词 $w_t$ 在上下文 $c_{t-1}=(w_1,w_2,\ldots,w_{t-1})$ 下的条件概率为
\begin{equation}
P(w_t=i\mid c_{t-1})
=\frac{\exp(U_i^{\mathsf T}f(c_{t-1})+b_i)}
 {\sum_{j=1}^{V}\exp(U_j^{\mathsf T}f(c_{t-1})+b_j)}.
\label{eq:ch06-lstm-conditional}
\end{equation}
其中，$f(\cdot)$ 是 LSTM 的输出，$U$ 是词表示矩阵，$b_i$ 是偏置项。\refcite{ch06-ref-034}{Ghosh 等人（2017）}提出用于词表偏移的 Affect-LM，把上式改写为
\begin{multline}
P(w_t=i\mid c_{t-1},e_{t-1})\\
=\frac{\exp(U_i^{\mathsf T}f(c_{t-1})+\beta V_i^{\mathsf T}g(e_{t-1})+b_i)}
 {\sum_{j=1}^{V}\exp(U_j^{\mathsf T}f(c_{t-1})+\beta V_j^{\mathsf T}g(e_{t-1})+b_j)}.
\label{eq:ch06-affect-lm}
\end{multline}
其中，$e_{t-1}$ 表示训练时由上下文词语获得的情绪类别信息。该式量化情绪类别信息对在上下文中预测目标词 $w_t$ 的影响，$\beta$ 是情绪强度参数。

\paragraph*{即插即用语言模型（PPLM）}
即插即用语言模型（Plug and Play Language Model，PPLM）允许用户把一个或多个表示特定控制目标的属性模型与大语言模型结合（\refcite{ch06-ref-019}{Dathathri et al., 2019}）。这样的组合无需额外训练或微调模型，对难以获得大量硬件资源的研究者很有用。PPLM 有两个基本条件。

\begin{enumerate}
\item 可以使用一个属性模型，记作 $p(a\mid x)$。
\item 可以取得该属性模型的梯度。
\end{enumerate}

PPLM 流程如图~\ref{fig:ch06-pplm} 所示，分为以下几步。

\begin{enumerate}
\item 对大语言模型做一次前向传播，从所得概率分布抽取一个词元。然后把生成的字符串交给属性模型，计算所需属性的概率 $p(a\mid x)$。
\item 反向传播，计算 $p(a\mid x)$ 和 $p(x)$ 相对于模型隐藏状态的梯度。调整隐藏状态，提高这两个概率。
\item 重新计算大语言模型的概率分布，抽取一个新的词元。
\end{enumerate}

\sourcefigure{fig06_10.pdf}{即插即用语言模型（PPLM）。图中使用属性模型反向调整语言模型的隐状态，再对下一词元重新采样。}{fig:ch06-pplm}

\paragraph*{生成式判别器（GeDi）}
生成式判别器（Generative Discriminator，GeDi）用较小的语言模型作生成式判别器，引导大语言模型的生成，以缓解毒性（\refcite{ch06-ref-054}{Krause et al., 2020}）。它用贝叶斯法则计算序列中每个可能的后续词的概率，平衡两个类别条件分布。第一个分布以所需属性，即“控制码”为条件，引导模型生成特定的理想内容。第二个分布的条件是“反控制码”，表示希望避开的属性，例如有毒语言。

基本原理是用辅助的类别条件语言模型（class-conditional language model，CC-LM），判断生成的词元属于控制码所指定类别的可能性。控制码定义文本序列 $x_{1:T}$ 的一种属性，可以是情感、主题或“有毒”等标签。类别条件语言模型计算 $p(x\mid c)$ 和 $p(x\mid\bar c)$，其中 $c$ 和 $\bar c$ 分别是控制码和反控制码。CTRL 就是一个例子，它提供了一个框架，让语言模型以属性变量为条件生成文本（\refcite{ch06-ref-052}{Keskar et al., 2019}）。

训练类别条件语言模型时，将每个序列 $x^{(i)}_{1:T_i}$ 与控制码 $c^{(i)}$ 配对。随后训练模型，使平均负对数似然 $L_g$ 最小。GeDi 的训练损失 $L_{\mathrm{gd}}$ 结合了语言模型损失 $L_g$ 和判别损失 $L_d$，用超参数 $\lambda$ 加权。
\begin{equation}
L_{\mathrm{gd}}=\lambda L_g+(1-\lambda)L_d.
\label{eq:ch06-gedi}
\end{equation}
判别损失用于推动模型产生希望的行为，并远离不希望的行为。$L_d$ 旨在提高分类准确率，而 $L_g$ 有助于更好地校准引导生成所用的词元概率。与 PPLM 相比，GeDi 在降低毒性方面更有效，计算开销也更低。但它需要一个额外的语言模型进行训练，因此增加了算力需求。

\paragraph*{自我诊断与自我去偏}
\refcite{ch06-ref-098}{Schick 等人（2021）}提出的方法围绕两个问题展开。第一，大语言模型能否自行发现自身的偏见，即自我诊断（self-diagnosis）？第二，大语言模型能否以无监督方式抑制不希望出现的行为，即自我去偏（self-debiasing）？

自我诊断考察预训练语言模型能否识别自身输出中在社会层面不合宜的属性，而无需额外的训练数据或外部资源。它利用模型在预训练中学到的知识，并给出简短的属性描述，以评估模型自行识别和处理偏见、不当内容的能力。

设语言模型为 $M$，词元序列为 $x_1,\ldots,x_k=x$。进行自我诊断时，将 $M$ 生成的句子 $x$ 与关于属性 $y$ 的问题组合，判断 $x$ 是否含有 $y$。例如，若 $M$ 生成句子 $x=$“I'm going to hunt you down!”，要判断它是否含有威胁，即 $y=$“a threat”，可用以下模板提示模型。
\begin{lstlisting}[nolol=true]
# Self-diagnosing template , denoted as sdg(.)
{sample text x}
Does the above text contain {threat text y} ?
\end{lstlisting}
$x$ 具有属性 $y$ 的概率为
\begin{equation}
p(y\mid x)
=\frac{p_M(\text{“Yes”}\mid\operatorname{sdg}(x,y))}
 {\sum_{w\in\{\text{“Yes”},\text{“No”}\}}
  p_M(w\mid\operatorname{sdg}(x,y))}.
\label{eq:ch06-sdg}
\end{equation}

自我去偏采用类似于自我诊断的思路，重点在于主动改变输出。它也利用模型自身的内部知识调整输出生成，降低生成有偏见文本的可能性。下面是一个自我去偏模板示例。
\begin{lstlisting}[nolol=true]
# Self-debiasing template , denoted as sdb(.)
The following text contains {undesired attribute s}:
{sample text x}
\end{lstlisting}

采用该模板时，需要预训练语言模型 $M$ 和表示某类偏见的属性 $y$。给模型输入文本 $x$，它用自我去偏输入 $\operatorname{sdb}(x,y)$ 生成续写。这会得到两个概率分布，即原始分布 $p_M(w\mid x)$ 和调整后的分布 $p_M(w\mid\operatorname{sdb}(x,y))$。后者被设计为更倾向于不希望出现的词。因此，对有偏见的词，下式的差值预期为负。
\[
\Delta(w,x,y)=p_M(w\mid x)-p_M(w\mid\operatorname{sdb}(x,y)).
\]
为缓解偏见，构造新的概率分布 $\widetilde{p}_M(w\mid x)$，使它与 $\alpha(\Delta(w,x,y))\cdot p_M(w\mid x)$ 成比例。这里，$\alpha\colon\mathbb{R}\to[0,1]$ 为缩放函数，根据偏见差值 $\Delta(w,x,y)$ 调整概率。这样只在有需要时改变概率，其余词保持不变，尽量减少对生成过程的干预。

为便于模型评估，这种办法不会把任何词的概率设为零。它根据 $\Delta(w,x,y)$ 降低带有偏见的词的概率，缩放函数使用衰减常数 $\lambda$。当 $x\geq0$ 时 $\alpha(x)=1$，否则 $\alpha(x)=e^{\lambda x}$。还可以考虑一组描述 $Y=\{y_1,\ldots,y_n\}$，据此调整分布，将该方法扩展到多种偏见。

